% 本文件由 03_代码/p4_tables.py 自动生成，请勿手工修改。
% 数据源：06_支撑材料/p4_results.json

\begin{table}[htbp]
  \centering
  \caption{附件 4 实时电价的结构分解（$P=F+R$，$F$ 为附件 1 的固定曲线，$R$ 为残差）}
  \label{tab:p4-struct}
  \small
  \begin{tabular}{lrr}
    \toprule
    项目 & 数值 & 说明 \\
    \midrule
    逐年逐段均值 $\overline{P}$ 与 $F$ 的最大偏差 & $5.15e-05$ & 两者\textbf{骨架与水平相同} \\
    $F$ 的年平均 & 0.7662 & 元/kWh \\
    $\overline{P}$ 的年平均 & 0.7662 & 元/kWh，与 $F$ 相差不足 $10^{-6}$ \\
    $R$ 与 $F$ 的相关系数 & $-0.0000$ & 正交，故 $R$ 就是``波动''本身 \\
    \midrule
    $R$ 的周一均值 & +0.0529 & 元/kWh \\
    $R$ 的周二均值 & +0.0549 & 元/kWh \\
    $R$ 的周三均值 & +0.0576 & 元/kWh \\
    $R$ 的周四均值 & +0.0542 & 元/kWh \\
    $R$ 的周五均值 & -0.1372 & 元/kWh \\
    $R$ 的周六均值 & -0.1364 & 元/kWh \\
    $R$ 的周日均值 & +0.0530 & 元/kWh \\
    \midrule
    周五、周六平均 & -0.1368 & 元/kWh，\textbf{偏低} \\
    其余五天平均 & +0.0545 & 元/kWh \\
    两者之差 & +0.1913 & 元/kWh，相当于均价的 25.0\% \\
    星期结构解释的 $R$ 方差 & 41.7\% & 其余为无规律的日内噪声 \\
    ``同星期加权''恒等式最大残差 & $0.00$ & 见式~\eqref{eq:p4-price} \\
    \bottomrule
  \end{tabular}
  \par\vspace{2pt}
  \begin{minipage}{\textwidth}\footnotesize
    注：附件 4 的电价可表示为附件 1 的 144 时段基准曲线与星期周期残差之和。两者年平均值相同，因此对照中的费用差异主要来自价格波动。周五、周六残差均值较低，据此采用同星期加权预测模型。
  \end{minipage}
\end{table}

\begin{table}[htbp]
  \centering
  \caption{三种价格信息条件下的策略在同一实际价格路径上的独立时序仿真（见第 6.1 节）}
  \label{tab:p4-strategy}
  \small
  \begin{tabular}{llrrrr}
    \toprule
    分支 & 策略 & 计划费/元 & 调整费/元 & 紧急费/元 & 合计费用/元 \\
    \midrule
    4-2 & 波动电价预测策略 & 14,041,884 & 0 & 698,929 & \textbf{14,740,813} \\
    4-2 & 固定电价参考策略 & 14,044,327 & 0 & 716,626 & \textbf{14,760,952} \\
    4-2 & 未来价格已知参考 & 13,911,363 & 0 & 735,304 & \textbf{14,646,668} \\
    \midrule
    4-2 & \textit{ΔJ（固定电价 − 波动电价）} & --- & --- & --- & \textbf{20,139}（+0.14\%） \\
    4-3 & 波动电价预测策略 & 13,802,989 & 197,228 & 153,069 & \textbf{14,153,286} \\
    4-3 & 固定电价参考策略 & 13,841,150 & 192,888 & 155,479 & \textbf{14,189,517} \\
    4-3 & 未来价格已知参考 & 13,726,007 & 176,274 & 173,391 & \textbf{14,075,672} \\
    \midrule
    4-3 & \textit{ΔJ（固定电价 − 波动电价）} & --- & --- & --- & \textbf{36,231}（+0.26\%） \\
    \bottomrule
  \end{tabular}
  \par\vspace{2pt}
  \begin{minipage}{\textwidth}\footnotesize
    注：三种策略均在同一条实际价格路径上结算。$\Delta J$ 定义为“固定电价参考”费用减去“波动电价预测”费用，用于衡量价格预测的经济效益。“未来价格已知参考”假设当天价格在 0:00 已全部公布，作为完全信息参照，不计入 $\Delta J$。三种策略分别依据各自的信息集独立标定。
  \end{minipage}
\end{table}

\begin{table}[htbp]
  \centering
  \caption{价格预测模型选择：在 2025 年 1 月预热期（评价区间之外）比较候选模型，并在子窗口上检验稳定性}
  \label{tab:p4-select}
  \small
  \begin{tabular}{lrrr}
    \toprule
    候选预测器 & 全窗口（30 天） & 后段（9 天） & 后半（16 天） \\
    \midrule
    \textbf{同星期加权 0.5/0.3/0.2} & \textbf{0.0486} & 0.0351 & 0.0341 \\
    同星期等权 1/3 & 0.0487 & 0.0353 & 0.0343 \\
    同星期加权 0.6/0.3/0.1 & 0.0488 & 0.0355 & 0.0344 \\
    同星期(上周同一天) & 0.0512 & --- & --- \\
    近 7 天滚动均值 & 0.0847 & --- & --- \\
    近 14 天滚动均值 & 0.0848 & --- & --- \\
    近 28 天滚动均值 & 0.0852 & --- & --- \\
    昨天同时段 & 0.0861 & --- & --- \\
    \midrule
    \multicolumn{4}{l}{\textit{选中的一组：权重 0.5/0.3/0.2，}\ $\beta=(1,\,0.75,\,0.5)$} \\
    \bottomrule
  \end{tabular}
  \par\vspace{2pt}
  \begin{minipage}{\textwidth}\footnotesize
    注：单位为元/kWh，指标为各候选在全天 $144$ 个时段的平均绝对误差。“全窗口”为第 $1$--$30$ 天；“后段”为滞后 $7/14/21$ 天样本均齐备的第 $22$--$30$ 天；“后半”为第 $15$--$30$ 天。三个窗口的最优候选一致，因此评价期内固定该模型。$\beta$ 按发布时刻分别选择，取值见上表最后一行。
  \end{minipage}
\end{table}

\begin{table}[htbp]
  \centering
  \caption{价格预测候选方案在 2025 年 2 月 1 日--12 月 31 日的平均绝对误差（\textbf{样本外复核}，不用于选型）}
  \label{tab:p4-forecast}
  \small
  \begin{tabular}{lrr}
    \toprule
    预测方案 & 平均绝对误差/(元/kWh) & 相对均价 \\
    \midrule
    同星期加权 0.6/0.3/0.1 & 0.0441 & 5.82\% \\
    \textbf{同星期加权 0.5/0.3/0.2（本文冻结使用）} & 0.0446 & 5.88\% \\
    同星期等权 1/3 & 0.0466 & 6.15\% \\
    同星期(上周同一天) & 0.0472 & 6.22\% \\
    近 7 天滚动均值 & 0.0827 & 10.91\% \\
    近 14 天滚动均值 & 0.0838 & 11.07\% \\
    昨天同时段 & 0.0843 & 11.13\% \\
    近 28 天滚动均值 & 0.0867 & 11.45\% \\
    全样本固定日内形状（事后口径，仅作下界） & 0.0964 & 12.72\% \\
    \bottomrule
  \end{tabular}
  \par\vspace{2pt}
  \begin{minipage}{\textwidth}\footnotesize
    注：“相对均价”以评价区间 334 天的均价 $0.7575$~元/kWh 为基准，全年均价为 $0.7662$~元/kWh。“全样本固定日内形状”使用评价区间均值，包含未来信息，仅作为误差参照。近 $N$ 天滚动均值未区分星期类型，因此未能提取星期偏移。评价区间内 $0.6/0.3/0.1$ 的误差为 $0.0441$，预先选定的 $0.5/0.3/0.2$ 为 $0.0446$，相差 $0.0005$ 元/kWh。
  \end{minipage}
\end{table}

\begin{table}[htbp]
  \centering
  \caption{日内对数偏差修正系数 $\beta$ 的扫描（对剩余时段的平均绝对误差，评价区间）}
  \label{tab:p4-beta}
  \small
  \begin{tabular}{lrrr}
    \toprule
    $\beta$ & 6:00 版 & 12:00 版 & 18:00 版 \\
    \midrule
    0 & 0.0488 & 0.0472 & 0.0428 \\
    0.25 & 0.0470 & 0.0448 & 0.0405 \\
    0.5 & 0.0461 & 0.0436 & \textbf{0.0403} \\
    0.75 & 0.0462 & \textbf{0.0436} & 0.0419 \\
    1 & \textbf{0.0470} & 0.0448 & 0.0447 \\
    1.25 & 0.0485 & 0.0468 & 0.0485 \\
    \bottomrule
  \end{tabular}
  \par\vspace{2pt}
  \begin{minipage}{\textwidth}\footnotesize
    注：$\beta=0$ 表示不作日内修正。加粗行 $\beta=(1,\,0.75,\,0.5)$ 由 1 月预热期选定，并在评价区间内固定；若在评价区间内选择，结果为 $\beta=0.5$。当 $\beta\ge1.25$ 时，修正幅度增大，预测误差随之上升。
  \end{minipage}
\end{table}

\begin{table}[htbp]
  \centering
  \caption{价格预测的四个发布版本在各自发布时间更新后的精度}
  \label{tab:p4-skill}
  \small
  \begin{tabular}{lrrrr}
    \toprule
    发布版本 & 覆盖时段 & 平均绝对误差 & 相对均价 & 共同时段误差 \\
    \midrule
    0点版本 & 0--144 & 0.0446 & 5.88\% & 0.0428 \\
    6点版本 & 36--144 & 0.0470 & 6.20\% & \textbf{0.0389} \\
    12点版本 & 72--144 & 0.0436 & 5.76\% & 0.0409 \\
    18点版本 & 108--144 & 0.0403 & 5.33\% & 0.0403 \\
    \midrule
    对照：全天用同一个均价 & 全天 & 0.2774 & 36.62\% & --- \\
    \bottomrule
  \end{tabular}
  \par\vspace{2pt}
  \begin{minipage}{\textwidth}\footnotesize
    注：“覆盖时段”为时段索引范围。各版本仅修正发布后的时段，因此平均绝对误差的原始统计区间不同。最后一列统一在 $18{:}00$--$24{:}00$ 上计算误差，用于比较四个版本；加粗者为该口径下的最小值。最后一行使用全天统一均价，不含日内形状信息，其误差高于四个分时预测版本。
  \end{minipage}
\end{table}

\begin{table}[htbp]
  \centering
  \caption{风险余量口径对照：以实际电价为权重的分位数 vs 等权分位数（正式区间按每 7 天抽样求均值）}
  \label{tab:p4-risk}
  \small
  \begin{tabular}{lrrrrr}
    \toprule
    $\alpha$ & 价格加权/(kWh) & 等权·线性/(kWh) & 等权·下确界/(kWh) & 比线性插值 & 比下确界 \\
    \midrule
    0.50 & 1.322 & -1.461 & -3.806 & --- & --- \\
    0.60 & 14.503 & 11.362 & 10.380 & +27.65\% & +39.71\% \\
    0.70 & 28.953 & 25.367 & 25.911 & +14.14\% & +11.74\% \\
    0.80 & 46.264 & 41.997 & 44.774 & +10.16\% & +3.33\% \\
    \bottomrule
  \end{tabular}
  \par\vspace{2pt}
  \begin{minipage}{\textwidth}\footnotesize
    注：“等权·线性”采用 \texttt{numpy} 的线性插值分位数；“等权·下确界”与“价格加权”均采用下确界定义，因此末列仅反映权重变化，倒数第二列还包含分位数定义差异。$\alpha=0.50$ 时余量为负，故相对变化以“---”表示。
  \end{minipage}
\end{table}

\begin{table}[htbp]
  \centering
  \caption{净负荷预测误差 $\varepsilon$ 与实际电价的协方差检验}
  \label{tab:p4-cov}
  \small
  \begin{tabular}{lr}
    \toprule
    项目 & 数值 \\
    \midrule
    样本数 & 48\,096 \\
    合并口径相关系数 & $+0.0863$ \\
    逐时段相关系数均值 & $+0.216$ \\
    逐时段相关系数范围 & $+0.085$~$+0.402$ \\
    逐时段相关为正的比例 & 100.0\% \\
    \midrule
    电价最高 20\% 时段的 $\varepsilon$ 均值 & $+40.5$~kW \\
    电价最低 20\% 时段的 $\varepsilon$ 均值 & $-74.7$~kW \\
    两者之差 & $+115.2$~kW \\
    \bottomrule
  \end{tabular}
  \par\vspace{2pt}
  \begin{minipage}{\textwidth}\footnotesize
    注：合并口径同时包含时段间价格差异；固定时段后，相关系数由 $+0.086$ 升至 $+0.216$，且全部时段均为正。因此，$\mathrm{Cov}(P,\varepsilon)>0$ 主要体现为同一时段的跨日相关性，风险余量据此按时段分别计算。
  \end{minipage}
\end{table}

\begin{table}[htbp]
  \centering
  \caption{波动电价下 4-3 的八种预报使用组合（各组合单独标定参数，评价区间与计费口径完全一致）}
  \label{tab:p4-combo}
  \small
  \begin{tabular}{lrrrrrr}
    \toprule
    使用的预报时刻 & 初始计划费/元 & 调整费/元 & 紧急费/元 & 合计费用/元 & 相对 $\varnothing$ & 期末储电量/kWh \\
    \midrule
    {6,12,18} & 13,802,989 & 197,228 & 153,069 & \textbf{14,153,286} & -5.79\% & 5,889.6 \\
    {6,18} & 13,802,526 & 302,659 & 153,716 & \textbf{14,258,902} & -5.09\% & 5,889.6 \\
    {12,18} & 13,880,759 & 83,222 & 374,151 & \textbf{14,338,131} & -4.56\% & 5,889.6 \\
    {6,12} & 14,076,549 & 16,306 & 277,579 & \textbf{14,370,434} & -4.34\% & 6,062.4 \\
    {18} & 13,880,487 & 236,489 & 386,855 & \textbf{14,503,831} & -3.46\% & 5,889.6 \\
    {12} & 14,242,219 & -77,058 & 412,319 & \textbf{14,577,480} & -2.96\% & 6,234.6 \\
    {6} & 13,974,367 & 82,343 & 548,690 & \textbf{14,605,400} & -2.78\% & 6,242.3 \\
    仅 0:00 & 14,209,197 & 0 & 813,704 & \textbf{15,022,901} & --- & 6,242.3 \\
    \bottomrule
  \end{tabular}
  \par\vspace{2pt}
  \begin{minipage}{\textwidth}\footnotesize
    注：与问题三的同一张表对照可以看出价格波动改变了什么——在固定电价下``用不用新预报''主要体现在调整费与紧急费的此消彼长；在波动电价下，更新预报同时改善了\textbf{电量}与\textbf{价格}两项判断（6:00 起既能看到当天更准的负载与光伏，也能用当天已实现的价格修正剩余时段的电价预测），因此增量收益的来源更宽。所有方案都使用 0:00 发布的预报，组合列出的时刻是额外使用的版本。
  \end{minipage}
\end{table}

\begin{table}[htbp]
  \centering
  \caption{两个分支在 334 个交付日上的检查结果（未通过天数均为零）}
  \label{tab:p4-check}
  \small
  \begin{tabular}{lrr}
    \toprule
    检查项 & result4-2 未通过天数 & result4-3 未通过天数 \\
    \midrule
    信息边界 & 0 & 0 \\
    计划冻结 & 0 & 0 \\
    物理可行性 & 0 & 0 \\
    状态累计 & 0 & 0 \\
    费用一致性 & 0 & 0 \\
    调整权限 & 0 & 0 \\
    \midrule
    \multicolumn{3}{l}{\textit{数值口径的最坏值}} \\
    \midrule
    逐段电量平衡最大残差/kWh & $0.00$ & $0.00$ \\
    计划费重算最大偏差/元 & $0.00$ & $0.00$ \\
    调整费重算最大偏差/元 & --- & $0.00$ \\
    紧急费重算最大偏差/元 & $0.00$ & $0.00$ \\
    跨日衔接最大偏差/kWh & $0.00$ & $0.00$ \\
    全局储电量范围/kWh & 1,200.0~10,800.0 & 1,200.0~10,800.0 \\
    \midrule
    储电下界余量/kWh & 0.000 & 0.000 \\
    储电上界余量/kWh & 0.000 & 0.000 \\
    \bottomrule
  \end{tabular}
  \par\vspace{2pt}
  \begin{minipage}{\textwidth}\footnotesize
    注：残差均小于 $10^{-12}$~kWh。储电量上下界余量为零，表明个别日期的容量约束达到边界。result4-2 的全天计划在 0:00 确定，执行阶段不发生第二次常规购电。
  \end{minipage}
\end{table}
